\documentclass[11pt]{article}
\usepackage{amsmath,amssymb,amsthm}
\usepackage[T1]{fontenc}
\usepackage{amsmath,amssymb,amsthm}
\usepackage[margin=1in]{geometry}

\newcommand{\Z}{\mathbb Z}
\newcommand{\Q}{\mathbb Q}
\newcommand{\NS}{\operatorname{NS}}
\newcommand{\rk}{\operatorname{rk}}
\newcommand{\Ocal}{\mathcal O}

\newtheorem{lemma}{Lemma}[section]
\newtheorem{proposition}[lemma]{Proposition}
\theoremstyle{remark}
\newtheorem{remark}[lemma]{Remark}

\newcommand{\fptitle}{Review of Problem 10 -- Astra}
\newcommand{\reviewtag}{Reviewer 1}
\usepackage{fp-review}

\begin{document}
\maketitle

\section{Review}

\paragraph{Correctness.}
\begin{itemize}
	\item  1: essentially flawless mathematical argument
	\item  0.5: minor errors that are easily corrected
	\item  0: blank solution, or acknowledged partial result
	\item  $-0.5$: major errors
	\item  $-1$: proposed proof makes no progress towards the result.
\end{itemize}
Note that points should be docked for wrong statements/ solutions.

\noindent Score: 0

The model is honest: "I cannot provide a complete proof." It checks some easy things, constituting perhaps 10\% towards a solution. 

\paragraph{Exposition.}
\begin{itemize}
	\item  1: well-written and clear solution.
	\item  0: argument that is possible to follow with serious effort
	\item  $-1$: very poor exposition
\end{itemize}

\noindent Score: 1

I think the elementary things that this model does write are explained clearly. 

\paragraph{Attributions.}
\begin{itemize}
	\item  1: solutions contains fully appropriate attributions
		to related work and theorems being used.
	\item  0: moderate issues with missing attributions or hallucinated references
	\item  $-1$: blatant plagiarism or hallucinations
\end{itemize}

\noindent Score: 1

I think the attributions for the elementary material included here (very little of it) are appropriate.

\newpage
\section{Submission with inline comments}


\begin{abstract}
I cannot provide a complete proof of the assertion in the question.
The argument below establishes its lattice-theoretic prerequisites and
identifies the remaining geometric existence problem.  It is not a proof
of the requested theorem.
\end{abstract}

\section{Scope and reference conventions}

The proposed conclusion specializes the forward existence assertion in
\cite[Conjecture~5.10]{GT24}, with the strengthened positivity hypothesis
in the question.
The fact that the general assertion is formulated as a conjecture does
not, by itself, establish that this specialization is open.
Nevertheless, citing that conjecture would not constitute a proof.

Here, references to numbered statements in \cite{GT24} concern
arXiv version~3, dated July~22, 2024.  Its Conjecture~5.10 is
Conjecture~5.12 in version~1.

Write
\[
 \Lambda=H^{\oplus 3}\oplus E_8^{\oplus 2},
 \qquad L=H,
 \qquad M=H\oplus E_8^{\oplus 2}.
\]
For a marked smooth fibre of the degeneration, let
\(I\subset\Lambda\) be its primitive rank-two isotropic monodromy
sublattice, as in \cite[Section~3.1]{GT24}.

Monodromy invariance gives \(I\subset L^\perp\).
For completeness, over \(\Q\), if \(N\) is the logarithm of monodromy,
then \(N\) is skew-adjoint for the intersection pairing and annihilates
\(L\).  Consequently,
\[
 (l,Nx)=-(Nl,x)=0
 \qquad(l\in L,\ x\in\Lambda_\Q).
\]
The same orthogonality holds for the saturation \(I\) of the image of
\(N\).

\section{The automatic lattice-theoretic conditions}

\begin{lemma}\label{lem:split}
Let \(T\) be an even unimodular lattice, and let
\(I\subset T\) be a primitive, totally isotropic sublattice of rank two.
There is an orthogonal decomposition
\[
 T\cong H_1\oplus H_2\oplus K
\]
in which \(I\) is generated by one primitive isotropic vector from each
of \(H_1\) and \(H_2\).
In particular, \(I\) is doubly admissible in the terminology of
\cite[Definition~3.31]{GT24}.
\end{lemma}

\begin{proof}
Choose a basis \(e_1,e_2\) of \(I\).
Since \(I\) is primitive, \(T/I\) is free.  Thus restriction gives a
surjection
\[
 \operatorname{Hom}(T,\Z)\longrightarrow
 \operatorname{Hom}(I,\Z).
\]
Unimodularity identifies \(T\) with \(\operatorname{Hom}(T,\Z)\).
We can therefore choose \(w_1,w_2\in T\) satisfying
\[
 (e_i,w_j)=\delta_{ij}.
\]
Since \(T\) is even, the vectors
\[
 g_i=w_i-\frac{w_i^2}{2}e_i
\]
are integral and isotropic, and still satisfy
\((e_i,g_j)=\delta_{ij}\).

Put \(c=(g_1,g_2)\), and replace \(g_2\) by \(g_2-ce_1\).
Its square remains zero because \((g_2,e_1)=0\).
After this replacement, \((g_1,g_2)=0\), while all the required
pairings with \(e_1,e_2\) remain unchanged.  Hence
\[
 U=\Z e_1+\Z g_1+\Z e_2+\Z g_2
\]
is isometric to \(H\oplus H\).

A unimodular sublattice splits off orthogonally: for each \(x\in T\),
unimodularity of \(U\) supplies \(u\in U\) with
\((u,v)=(x,v)\) for every \(v\in U\), and then \(x-u\in U^\perp\).
Thus \(T=U\oplus U^\perp\), as required.
\end{proof}

\begin{proposition}\label{prop:degeneration}
For the degeneration in the question, there is a negative-definite
even unimodular lattice \(K\) of rank sixteen such that
\[
 L^\perp\cong H^{\oplus2}\oplus K,
 \qquad
 I^\perp_\Lambda/I\cong L\oplus K.
\]
Under the second isomorphism, the polarization is the indicated
summand \(L\).  In particular, its embedding in \(I^\perp_\Lambda/I\)
is primitive, and the degeneration is doubly admissible.
\end{proposition}

\begin{proof}
Because \(L=H\) is unimodular, its embedding in \(\Lambda\) splits
orthogonally:
\[
 \Lambda=L\oplus T,
 \qquad T=L^\perp.
\]
The lattice \(T\) is even unimodular of signature \((2,18)\).
Apply Lemma~\ref{lem:split} to \(I\subset T\).
The remaining summand \(K\) has signature \((0,16)\).
In the resulting decomposition, the orthogonal complement of \(I\)
inside \(H_1\oplus H_2\) is exactly \(I\).  Taking the quotient by
\(I\) therefore gives
\[
 I^\perp_\Lambda/I=L\oplus K.
\]
All the assertions follow.
\end{proof}

\begin{lemma}\label{lem:mirror}
Every primitive vector \(e\in I\) is \(1\)-admissible in \(L^\perp\).
Its mirror lattice is
\[
 e^\perp_{L^\perp}/\Z e\cong M.
\]
\end{lemma}

\begin{proof}
Set \(T=L^\perp\).
Since \(e\) is primitive and \(T\) is unimodular, there is a vector
\(w\in T\) with \((e,w)=1\).
The vector
\[
 g=w-\frac{w^2}{2}e
\]
is integral and isotropic, with \((e,g)=1\).
Thus \(\Z e+\Z g\cong H\) splits off from \(T\).
This proves \(1\)-admissibility and shows that
\(e^\perp_T/\Z e\) is even unimodular of signature \((1,17)\).
The classification of indefinite even unimodular lattices identifies
it with \(H\oplus E_8^{\oplus2}\); see, for example,
\cite[Section~1.14]{Nik80}.
\end{proof}

\section{An elliptic fibration with the required abstract frame lattice}

The preceding lattice calculation also produces an elliptic fibration
on \(Y\) whose abstract frame lattice is the correct one.
It does not produce a splitting loop.

\begin{proposition}\label{prop:fibration}
Let \(Y\) be as in the question, and let \(K\) be the lattice from
Proposition~\ref{prop:degeneration}.
There is an elliptic fibration
\[
 f:Y\longrightarrow\mathbb P^1
\]
with a section and fibre class \(F\), such that
\[
 F^\perp_{\NS(Y)}/\Z F\cong K.
\]
\end{proposition}

\begin{proof}
Very generality gives \(\NS(Y)=M\).
The lattice \(H\oplus K\) is even unimodular of signature \((1,17)\);
hence it is isometric to \(M\).

Choose a primitive isotropic generator of the factor \(H\) in this
decomposition of \(\NS(Y)\).
After changing its sign and applying reflections in classes of
\((-2)\)-curves, it becomes a primitive nef class \(F\).
The standard reduction to the nef cone and the elliptic-pencil
theorem for primitive nef isotropic classes give an elliptic
fibration with fibre class \(F\); see
\cite[Chapter~8, Section~2, and Chapter~11, Section~1]{Huy16}.
Since these reflections are lattice isometries,
\[
 F^\perp_{\NS(Y)}/\Z F\cong K.
\]

We check the existence of a section.
Unimodularity supplies \(w\in\NS(Y)\) with \((F,w)=1\).
Replacing \(w\) by \(w-\frac{w^2}{2}F\), we may suppose that \(w^2=0\).
Then
\[
 D=w-F
 \quad\text{satisfies}\quad
 D^2=-2,\qquad D\cdot F=1.
\]
Riemann--Roch on a K3 surface implies that \(D\) or \(-D\) is effective.
The latter is impossible because \(F\) is nef and
\((-D)\cdot F=-1\).
An effective representative of \(D\) therefore has a horizontal
irreducible component of degree one over \(\mathbb P^1\).
Its normalization gives a section.
\end{proof}

This is the coarse degeneration--fibration correspondence: it matches
the relevant abstract lattices, but not the componentwise splitting
required in the question.  Compare \cite[Proposition~4.1]{DHT17}.

\section{The missing geometric assertion}

Suppose that an allowable loop \(\gamma\) for the fibration in
Proposition~\ref{prop:fibration} has been constructed and that its
associated class \(\tau\) lies in \(\NS(Y)^\perp\).
Put
\[
 Q_X=I^\perp_\Lambda/I,\qquad
 Q_Y=(\Z F\oplus\Z\tau)^\perp/(\Z F\oplus\Z\tau),\qquad
 \Gamma=F^\perp_{\NS(Y)}/\Z F.
\]
Use the component sublattices from \cite[Definition~5.1]{GT24}:
\[
 P_i^X=\varphi_X(K_{V_i}^{\perp}),
 \qquad
 P_i^Y=\varphi_Y([r_i(b)]^{\perp}).
\]
Here \(r_i\) is the right adjoint of the asymptotic charge map of
the fibration over the \(i\)-th disc.

The two polarizations are then doubly admissible.
Indeed, this has already been proved for the degeneration.
On the fibration side, the section makes \(F\) \(1\)-admissible.
Moreover, \(\NS(Y)^\perp\cong H^{\oplus2}\) is unimodular, so the
primitive isotropic class \(\tau\) is \(1\)-admissible by the argument
of Lemma~\ref{lem:mirror}.

Consequently, \cite[Proposition~5.3]{GT24} reduces the mirror-pair
condition to the existence of an isometry
\[
 \psi:Q_Y\xrightarrow{\sim}Q_X
 \quad\text{such that}\quad
 \psi(\Gamma)=L^\perp_{Q_X},
 \qquad
 \psi(P_i^Y)=P_i^X\quad(i=1,2).
 \tag{*}
\]

The unresolved requirements are therefore:
\begin{enumerate}
 \item construct an allowable loop compatible with the polarization;
 \item verify the componentwise equalities in \((*)\) for the
       \emph{given} Tyurin model.
\end{enumerate}
Allowability includes both a specified boundary-monodromy condition
and a primitivity condition; see
\cite[Assumption~4.9 and Definition~4.11]{GT24}.

Neither requirement follows merely from the abstract isomorphism
\(\Gamma\cong K\).
In particular, an isometry of the ambient lattices need not identify
the two distinguished component sublattices.

The Type~II constructions in \cite[Section~7F]{ABE22} provide standard
divisor models.
Using them here would require an additional comparison with the
given polarized Tyurin model, followed by the loop construction and
the integral componentwise verification above.

\section*{Conclusion}

The calculations above verify the mirror lattice, double
admissibility, and the existence of an elliptic fibration with the
correct abstract frame lattice.  The sources also provide the precise
conditional criterion for being a mirror pair.

I have not established the required allowable-loop construction and
componentwise matching for every degeneration satisfying the
strengthened hypotheses.  Thus a complete proof has not been
achieved, and presenting the preceding reduction as one would leave
the central geometric assertion unproved.

\begin{thebibliography}{ABE22}

\bibitem[ABE22]{ABE22}
V.~Alexeev, A.~Brunyate, and P.~Engel,
\emph{Compactifications of moduli of elliptic K3 surfaces:
Stable pair and toroidal},
Geom.\ Topol.\ \textbf{26} (2022), no.~8, 3525--3588.

\bibitem[DHT17]{DHT17}
C.~F.~Doran, A.~Harder, and A.~Thompson,
\emph{Mirror symmetry, Tyurin degenerations and fibrations on
Calabi--Yau manifolds},
String-Math 2015, Proc.\ Sympos.\ Pure Math., vol.~96,
American Mathematical Society, 2017, pp.~93--131.

\bibitem[GT24]{GT24}
L.~Giovenzana and A.~Thompson,
\emph{Degenerations and fibrations of K3 surfaces:
Lattice polarisations and mirror symmetry},
arXiv:2405.12009v3, July~22, 2024.

\bibitem[Huy16]{Huy16}
D.~Huybrechts,
\emph{Lectures on K3 Surfaces},
Cambridge Studies in Advanced Mathematics, vol.~158,
Cambridge University Press, 2016.

\bibitem[Nik80]{Nik80}
V.~V.~Nikulin,
\emph{Integral symmetric bilinear forms and some of their applications},
Math.\ USSR Izv.\ \textbf{14} (1980), no.~1, 103--167.

\end{thebibliography}

\end{document}
